
\documentclass[aip,apl,reprint]{revtex4-2}
\begin{document}
\title{AIP Family: Official AIP Submission Preview}
\author{A. Placeholder}
\affiliation{Department of Verified Templates, Placeholder Institute, Placeholder City, Country}
\email{verified@example.org}
\author{B. Example}
\affiliation{Journal Systems Lab, Example University, Example City, Country}
\begin{abstract}
This submission-style preview is rendered with the official REVTeX class in AIP journal mode. The manuscript remains generic in scientific content, but the page is structured to expose frontmatter, section hierarchy, figure flow, table treatment, and bibliography rhythm in a form closer to a realistic AIP submission.
\end{abstract}
\maketitle
\section{Introduction}
The introduction is written as concise technical prose so the preview feels like a practical AIP manuscript rather than a minimal template test page.
\section{Related Work}
The related-work section is explicit so the section hierarchy is easier to inspect when the page transitions from context to method and experiment.
\section{Methodology}
The methodology section introduces representative mathematics and a figure placeholder so the REVTeX layout can be reviewed under more realistic technical material.
\begin{equation}
L = \sum_{i=1}^{M} w_i \left\| y_i - \hat{y}_i \right\|_2^2.
\end{equation}
\begin{figure}[t]
\centering
\fbox{\parbox[c][1.9in][c]{0.84\linewidth}{\centering AIP Figure Placeholder\\[0.4em]\normalsize Optical setup, measurement layout, or device schematic}}
\caption{Representative AIP figure used to inspect REVTeX spacing, caption placement, and the transition from methods text to the first major figure.}
\end{figure}
\section{Experiments}
The experiments section provides enough content to expose quantitative comparison rhythm and table placement inside the AIP submission layout.
\subsection{Experimental Setup}
This paragraph stands in for sample preparation, hardware settings, measurement conditions, and baseline methods.
\subsection{Quantitative Results}
\begin{table}[t]
\caption{Quantitative comparison block used to inspect the official AIP submission template under realistic manuscript content.}
\centering
\begin{tabular}{lcc}
\hline
Method & Score & Time \\
\hline
Placeholder A & 0.91 & 1.4 \\
Placeholder B & 0.88 & 1.7 \\
Placeholder C & 0.84 & 2.0 \\
\hline
\end{tabular}
\end{table}
\section{Conclusion and Discussion}
This page verifies the official AIP template path while preserving the six-part manuscript skeleton requested for the library. The preview is intended to read more like a review-format physics submission than a bare placeholder.
\begin{thebibliography}{9}
\bibitem{ref1} A. Smith and B. Lee, ``Template-aware AIP manuscript preparation,'' \textit{AIP Placeholder Journal} \textbf{1}, 1--10 (2025).
\bibitem{ref2} C. Garcia and D. Patel, ``Figure and table balance in REVTeX submissions,'' \textit{Appl. Phys. Lett.} \textbf{124}, 123456 (2026).
\end{thebibliography}
\end{document}
